\usepackage[T1]{fontenc}
\usepackage[utf8]{inputenc}
\usepackage{lmodern,microtype,amsmath,amssymb,amsthm,mathtools,thmtools,bm}
\usepackage[hmargin=0.94in,vmargin=0.82in]{geometry}
\usepackage{float}
\usepackage{booktabs,array,tabularx,longtable,enumitem,siunitx,xcolor,needspace}
\usepackage{tikz}
\usetikzlibrary{matrix,arrows.meta,positioning,calc,fit,backgrounds,decorations.pathreplacing}
\usepackage{pgfplots}
\pgfplotsset{compat=1.18}
\usepackage[ruled,vlined,linesnumbered]{algorithm2e}
\usepackage[numbers,sort&compress]{natbib}
\usepackage{xr}
\usepackage[hidelinks]{hyperref}
\usepackage[nameinlink,noabbrev]{cleveref}
\hypersetup{pdftitle={Distribution-independent SQ learning does not imply low dimension complexity},pdfauthor={Samuel Mausberg},pdfsubject={A negative answer to Open Question 2 of Feldman, Kamath and Srebro},pdfkeywords={statistical queries, sign-rank, dimension complexity, learning theory}}
\declaretheorem[name=Theorem]{theorem}
\declaretheorem[name=Lemma,sibling=theorem]{lemma}
\declaretheorem[name=Proposition,sibling=theorem]{proposition}
\declaretheorem[name=Corollary,sibling=theorem]{corollary}
\declaretheorem[name=Main Theorem,numbered=no]{maintheorem}
\declaretheorem[name=Auxiliary Lemma,numbered=no]{auxlemma}
\declaretheorem[name=Open Problem,numbered=no]{openproblem}
\crefname{theorem}{Theorem}{Theorems}
\crefname{lemma}{Lemma}{Lemmas}
\crefname{proposition}{Proposition}{Propositions}
\crefname{corollary}{Corollary}{Corollaries}
\crefname{algocf}{Algorithm}{Algorithms}
\numberwithin{equation}{section}
\DeclareMathOperator{\dc}{dc}
\DeclareMathOperator{\sr}{signrank}
\DeclareMathOperator{\sign}{sign}
\DeclareMathOperator{\tr}{tr}
\DeclareMathOperator{\diag}{diag}
\DeclareMathOperator{\rect}{rect}
\DeclareMathOperator{\RSD}{RSD}
\newcommand{\R}{\mathbb R}
\newcommand{\E}{\mathbb E}
\newcommand{\1}{\mathbf 1}
\newcommand{\cH}{\mathcal H}
\newcommand{\op}{\mathrm{op}}
\definecolor{oiBlue}{HTML}{0072B2}
\definecolor{oiOrange}{HTML}{E69F00}
\definecolor{oiGreen}{HTML}{009E73}
\definecolor{oiPurple}{HTML}{CC79A7}
\definecolor{oiGray}{HTML}{555555}
\tikzset{>=Latex,diagram box/.style={draw=oiGray,rounded corners=2pt,align=center,inner sep=5pt,font=\small},solid implication/.style={->,line width=.7pt},coexist/.style={dashed,line width=.7pt,draw=oiPurple}}
\SetAlFnt{\small}
\SetAlCapFnt{\small\bfseries}
\SetAlgoNlRelativeSize{-1}
\SetKw{Return}{return}
\SetKw{Continue}{continue}
\setlength{\emergencystretch}{2em}
\setlength{\textfloatsep}{14pt plus 3pt minus 2pt}
\setlength{\intextsep}{12pt plus 2pt minus 2pt}
\setlist[itemize]{topsep=4pt,itemsep=2pt}
\setlist[enumerate]{topsep=4pt,itemsep=2pt}
\renewcommand{\topfraction}{.93}
\renewcommand{\bottomfraction}{.85}
\renewcommand{\textfraction}{.06}
\renewcommand{\floatpagefraction}{.75}
\setcounter{topnumber}{3}
\setcounter{bottomnumber}{3}
\widowpenalty=7000
\clubpenalty=7000
\raggedbottom

\renewcommand{\theHequation}{\theHsection.\arabic{equation}}
\renewcommand{\theHtheorem}{\theHsection.\arabic{theorem}}
\providecommand{\theHlemma}{}
\renewcommand{\theHlemma}{\theHsection.\arabic{theorem}}
\providecommand{\theHproposition}{}
\renewcommand{\theHproposition}{\theHsection.\arabic{theorem}}
\providecommand{\theHcorollary}{}
\renewcommand{\theHcorollary}{\theHsection.\arabic{theorem}}

\crefname{appendix}{Appendix}{Appendices}
\Crefname{appendix}{Appendix}{Appendices}

\crefname{section}{Section}{Sections}
\Crefname{section}{Section}{Sections}
\crefname{subsection}{Section}{Sections}
\Crefname{subsection}{Section}{Sections}
